\chapter{For plugin developers}
\label{ch:developers}

\fxme\ runs any VST3 effect. A plugin can also be made \emph{\fxme-aware}, and
then it gets tier \tier{2} (chapter~\ref{ch:tiers}): its editor on the
client shows the live meters, spectra and values of the copy that runs on
the server. This chapter summarises what that takes. The reference is the
file \file{fxmenet-howto.md} of FxmeTools, with the demonstration plugin
\file{Examples/GainPan} in the \fxme\ source.

\section{The idea}

The plugin is loaded twice: a \emph{source} on the server, which processes
audio, and a \emph{mirror} on the client, which never processes and shows
the editor. What the editor reads that is not a parameter (a meter, a
spectrum, a computed value) is sent from the source to the mirror as
telemetry and written into the mirror's own objects, so the editor itself
does not change. Outside \fxme\ the plugin behaves exactly as before.

\section{What to do}

\begin{enumerate}
  \item Add an \texttt{fxme::net::NetBridge} to the processor, and return it from
    \texttt{getVST3Client\allowbreak Extensions()}.
  \item Register everything the editor reads that is not a parameter, with a
    stable name each: meters, taps, atomics, computed scalars.
  \item In the few getters that compute from the DSP, return the telemetry
    value when the plugin is a mirror.
  \item Skip the work that only feeds the DSP when the plugin is a mirror.
  \item Call \texttt{notifyStateChanged()} when state that is not a parameter
    changes (a loaded file, for example).
\end{enumerate}

\begin{lstlisting}[style=fxcode]
#include <FxmeTools/net/NetBridge.h>

class MyProcessor : public juce::AudioProcessor
{
public:
    juce::VST3ClientExtensions* getVST3ClientExtensions() override { return &net; }
private:
    Compressor compressor;
    fxme::net::NetBridge net;   // after everything it registers
};

MyProcessor::MyProcessor()
{
    net.telemetry().addMeter ("gr", compressor.getGrMeter());
}
\end{lstlisting}

\begin{longtable}{@{}p{0.34\textwidth}p{0.6\textwidth}@{}}
\toprule
\textbf{The editor reads} & \textbf{Register}\\
\midrule
\endhead
a level meter & \texttt{addMeter (name, meter)}\\
a spectrum or waveform tap & \texttt{addTap (name, tap)}\\
a \texttt{std::atomic<float>} & \texttt{addValue (name, atomic)}\\
a function that computes a value & \texttt{addScalar (name, getter)}\\
a read-and-reset function & \texttt{addPeak (name, function)}\\
\bottomrule
\end{longtable}

\section{Checking}

The \fxme\ source includes a test tool, \cmd{FxmeNetLoopback}, which loads
your plugin twice in one process, as source and mirror, and pumps the
telemetry between them. Both editors appear side by side; the mirror's
meters must follow the source's.

\begin{lstlisting}
FxmeNetLoopback path/to/MyPlugin.vst3
\end{lstlisting}

\section{The \fxme\ compatible mark}
\label{sec:stamp}

Plugins that support \fxme\ may carry the \fxme\ mark, so that users know the
plugin gives a tier \tier{2} experience.

\begin{center}
  \fxlogo[width=0.28\textwidth]{FxmeNetLogo3.jpg}\qquad
  \fxlogo[width=0.08\textwidth]{fxmenet-logo.png}
\end{center}

\todo{The two marks above are drafts (a connector-shaped logo and a small
icon). Decide the final design and supply it as SVG and as transparent PNG in
several sizes. Write the conditions of use: what a plugin must do to carry it
(for example, that it passes the Loopback check and that its editor works
as a mirror), where to show it (website, manual, plugin's information
panel), and the licence of the artwork.}
